\chapter{Вся машина}
\chaptermark{Вся машина}
\label{sec:synthesis}

\section{Что мы собрали}

Мы брали по одному типу нейронов из таблицы~\ref{tab:types} и спрашивали, что он добавляет к вычислительной машине. Каждый раз правила были одни и те же: только то, что бывает в живом организме, и никаких часов. Теперь соберём части в одну машину и посмотрим, как они работают вместе.

Картина из главы~\ref{sec:neurontypes} подтвердилась и стала точнее. Огромная адресная сеть \nt{ГЛУТ}-нейронов несёт данные. \nt{ГАМК}-нейроны управляют потоком этих данных. А над сетью висят четыре широковещательных канала, каждый со своей ручкой: \nt{ДОФА} говорит, какие изменения весов оставить, \nt{СЕРО} держит общий уровень и, по гипотезе, горизонт планирования, \nt{НОРА} выбирает между поиском и решением, \nt{АЦЕТ} --- между записью и чтением (рис.~\ref{fig:machine}).

\begin{figure}[t]
\centering
\begin{tikzpicture}[>=Stealth,
  blk/.style={draw,very thick,rounded corners=3pt,align=center,font=\small},
  reg/.style={draw,very thick,double,double distance=1.2pt,circle,minimum size=1.25cm,inner sep=0pt,font=\scriptsize,fill=orange!15},
  bc/.style={->,thick,densely dashed,orange!85!black}]
% sensors, bus, actions
\node[blk,fill=green!8,text width=1.7cm] (S) at (0.9,0) {датчики\\\scriptsize вкл./выкл.};
\draw[very thick,rounded corners=4pt,fill=blue!5] (2.6,-1.35) rectangle (9.0,1.35);
\node[font=\small] at (5.8,0.95) {шина \nt{ГЛУТ}: данные};
\node[font=\scriptsize,align=center] at (4.3,0.25) {совпадения, задержки,\\логика (гл.~\ref{sec:glutcomputer})};
\node[font=\scriptsize,align=center] at (7.4,0.25) {память, код\\(гл.~\ref{sec:code}, \ref{sec:acet})};
\draw[thick,red!70!black,fill=red!8,rounded corners=2pt] (2.8,-1.15) rectangle (8.8,-0.55);
\node[font=\scriptsize,red!60!black] at (5.8,-0.85) {\nt{ГАМК} (гл.~\ref{sec:gaba}): запрет, выбор, деление, ритм};
\node[blk,fill=blue!5,text width=2.0cm] (A) at (10.9,0) {выбор\\действия};
\draw[->,very thick] (S) -- (2.6,0);
\draw[->,very thick] (9.0,0) -- (A);
\draw[->,very thick] (A) -- (12.6,0) node[right,font=\small] {мышцы};
\pic[scale=1.1] at (13.2,-0.5) {spring};
\pic[scale=1.15] at (0.4,0.85) {eyepic};
\pic[scale=1.05] at (1.35,0.85) {speaker};
% registers
\node[reg] (D) at (3.4,3.2) {\nt{ДОФА}};
\node[reg] (C) at (5.6,3.2) {\nt{СЕРО}};
\node[reg] (N) at (7.8,3.2) {\nt{НОРА}};
\node[reg] (H) at (10.0,3.2) {\nt{АЦЕТ}};
\node[font=\scriptsize,align=center,above=1pt] at (D.north) {ошибка $\delta$};
\node[font=\scriptsize,align=center,above=1pt] at (C.north) {уровень, $\tau_\gamma$};
\node[font=\scriptsize,align=center,above=1pt] at (N.north) {усиление $g$};
\node[font=\scriptsize,align=center,above=1pt] at (H.north) {запись $a$};
\draw[bc] (D.south) -- (3.4,1.35);
\draw[bc] (C.south) -- (5.6,1.35);
\draw[bc] (N.south) -- (7.8,1.35);
\draw[bc] (H.south) -- (10.0,2.1) -- (8.8,1.35);
\draw[bc] (H.south east) to[bend left=20] (A.north east);
\draw[->,thick] (2.85,1.35) -- (2.85,2.75) -- (D.west) node[pos=0.25,left,font=\scriptsize] {$V$};
\draw[->,thick,green!45!black] (0.4,3.2) node[left,black,font=\small] {награда $r$} -- (2.2,3.2) -- (D.west);
\pic[scale=0.9] at (-0.45,3.7) {juice};
\draw[->,thick,densely dotted,gray!50!black] ([yshift=0.45cm]D.north) to[bend left=18] node[above,font=\scriptsize,black] {неожиданность} ([yshift=0.45cm]N.north);
\end{tikzpicture}
\caption{Машина целиком. Шина \nt{ГЛУТ} несёт данные от датчиков к выбору действия; \nt{ГАМК} управляет потоком внутри неё. Двойные кружки --- широковещательные каналы; штриховые стрелки --- рассылка их сигналов всей сети, у каждого своя ручка. \nt{ДОФА} получает награду $r$ и ожидание $V$ от критика внутри шины (гл.~\ref{sec:dofa}); точечная стрелка --- гипотеза, что о неожиданности \nt{НОРА} узнаёт по необычно большим ошибкам (гл.~\ref{sec:nora}).}
\label{fig:machine}
\end{figure}

\section{Одна формула для всех ручек}

Ручки можно собрать в одну схематическую формулу. Это не новая модель, а сводка моделей глав~\ref{sec:glutcomputer}--\ref{sec:acet}: каждый символ взят из своей главы.
\begin{empheq}[box=\keybox]{equation}
\begin{aligned}
\tau\,\frac{\dd r_i}{\dd t} \;&=\; -r_i + F\Bigl(g\,\Bigl[\tfrac{1+a}{2}\,x_i + (1-a)\sum_j W_{ij}\,r_j - \sum_k M_{ik}\,q_k - \theta\Bigr]\Bigr),\\
\frac{\dd W_{ij}}{\dd t} \;&=\; \eta\,a\,\delta(t)\,e_{ij}(t),
\qquad \tau_e\,\frac{\dd e_{ij}}{\dd t} = -e_{ij} + r_i\,r_j,
\qquad \delta = r + \frac{\dd V}{\dd t} - \frac{V}{\tau_\gamma} .
\end{aligned}
\label{eq:machine}
\end{empheq}
Здесь $r_i$ --- частоты \nt{ГЛУТ}-нейронов, $x_i$ --- вход от датчиков, $W_{ij}\ge0$ --- их взаимные веса (глава~\ref{sec:glutcomputer}); $q_k$ --- частоты \nt{ГАМК}-нейронов, $M_{ik}\ge0$ --- сила их торможения (глава~\ref{sec:gaba}); $e_{ij}$ --- след, $\delta$ --- ошибка \nt{ДОФА} (глава~\ref{sec:dofa}); $\tau_\gamma$ --- горизонт, по гипотезе задаваемый \nt{СЕРО}; $g$ --- усиление \nt{НОРА} (глава~\ref{sec:nora}); $a$ --- уровень \nt{АЦЕТ} (глава~\ref{sec:acet}).

Посмотрите, чего в~\eqref{eq:machine} нет. Нет такта: всё записано дифференциальными уравнениями, и каждый нейрон меняется сам. Нет отдельной памяти: помнят те же $W_{ij}$, по которым идёт вычисление, и та же активность $r_i$. Нет персональных поправок для огромного большинства синапсов: их обучение --- произведение местного следа и одного общего числа; отдельный провод ошибки к каждому выходу есть лишь в цепи, которая учит движения (глава~\ref{sec:motorlearning}). И управляющих сигналов всего четыре рода: $\delta$, $\tau_\gamma$, $g$, $a$, на восемьдесят миллиардов нейронов. Каждый из них на самом деле не одно число, а небольшой жгут параллельных линий (утверждение~\ref{prop:bundle}), но линий в жгуте единицы.

\begin{keyblock}
\begin{proposition}[Архитектура]
\label{prop:architecture}
Машину мозга, насколько мы её сейчас понимаем, можно описать тремя слоями. Данные идут по адресной шине \nt{ГЛУТ}. Поток данных направляют, ограничивают и нормируют адресные \nt{ГАМК}-нейроны. Режим работы задают несколько широковещательных каналов, каждый --- небольшой жгут параллельных линий, общих для больших участков сети. Первые два слоя установлены твёрдо; распределение ролей между широковещательными каналами --- в основном гипотеза.
\end{proposition}
\end{keyblock}

\section{Все типы в одной таблице}

Таблица~\ref{tab:synthesis} сводит все типы нейронов книги. Четыре типа, которым не досталось своей главы, заняли в ней по строчке; двое из них нашли роль в главах о выходах машины: \nt{ГЛИЦ} --- в петлях спинного мозга (глава~\ref{sec:muscles}), \nt{АДРЕ} --- в химическом выходе (глава~\ref{sec:chemoutput}). Роль остальных двух в вычислении либо простая, либо неизвестна. Вещества, которые не задают тип нейрона, собраны в приложении~\ref{sec:othermediators}.

\begin{table}[t]
\centering
\footnotesize
\setlength{\tabcolsep}{3pt}
\renewcommand{\arraystretch}{1.15}
\begin{tabular}{>{\raggedright}p{1.3cm}>{\raggedright}p{1.0cm}>{\raggedright}p{3.9cm}>{\raggedright}p{1.5cm}>{\raggedright}p{1.8cm}>{\raggedright\arraybackslash}p{3.8cm}}
\toprule
Тип & Глава & Что вычисляет & Время & Шина & Что известно \\
\midrule
\nt{ГЛУТ} & \ref{sec:glutcomputer}, \ref{sec:code}, \ref{sec:sensors}, \ref{sec:motorlearning}, \ref{sec:dendrites} & данные: совпадения, задержки, логика, память, последовательности & мс & адресная & установлено \\
\nt{ГАМК} & \ref{sec:gaba}, \ref{sec:code}, \ref{sec:pain}, \ref{sec:motorlearning} & управление потоком: запрет, выбор, деление, устойчивость, ритм; ворота боли & мс & адресная & установлено; деление частоты --- вместе с фоновым шумом \\
\nt{СЕРО} & \ref{sec:serocomputer}, \ref{sec:pain}, \ref{sec:sleep} & регулятор уровня, сглаживание; горизонт $\tau_\gamma$; громкость боли; режимы сна & секунды & объём & самоторможение измерено; горизонт --- гипотеза \\
\nt{ДОФА} & \ref{sec:dofa}, \ref{sec:dofaalt} & ошибка предсказания $\delta$; медленный уровень --- цена времени & $0.1$\,с и минуты & широковещ. & $\delta$ измерено; расширения --- гл.~\ref{sec:dofaalt} \\
\nt{НОРА} & \ref{sec:nora}, \ref{sec:pain}, \ref{sec:chemoutput}, \ref{sec:sleep} & усиление $g$: искать или решать; неожиданность; «газ» для органов & секунды & широковещ. & рост отношения сигнала к шуму измерен; роль в поиске --- гипотеза \\
\nt{АЦЕТ} & \ref{sec:acet}, \ref{sec:muscles}, \ref{sec:chemoutput}, \ref{sec:sleep} & запись или чтение; скорость обучения; выход на мышцу; «тормоз» для органов & секунды & обе & три эффекта измерены; переключение режимов --- гипотеза \\
\nt{ГЛИЦ} & \ref{sec:muscles}, \ref{sec:pain} & отрицательный вес в спинном мозге и стволе: запрет мышцы-противника & мс & адресная & установлено \\
\nt{ГИСТ} & --- & глобальный сигнал «бодрствуй» & медленно & широковещ. & роль в вычислении не изучена \\
\nt{АДРЕ} & \ref{sec:chemoutput} & «газ» для всего тела через кровь; в мозге неизвестно & медленно & широковещ. & вне мозга установлено; роль в мозге неясна \\
\nt{АСПА} & --- & слабый положительный вес & мс & адресная & роль спорна \\
\bottomrule
\end{tabular}
\caption{Все типы нейронов книги: что вычисляет каждый и насколько это известно.}
\label{tab:synthesis}
\end{table}

\section{Как ручки работают вместе}

До сих пор каждая глава поворачивала одну ручку. Проследим теперь одно событие через всю машину (рис.~\ref{fig:timeline}). Животное давно выучило, что действие $A$ приносит сок. В какой-то момент мир меняется: $A$ перестаёт приносить сок, а приносит его действие $B$, которое животное почти не пробует.

\emph{Ошибка.} Сок после $A$ не пришёл, и \nt{ДОФА}-нейроны отвечают провалом, $\delta<0$ по формуле ошибки~\eqref{eq:ctd}. Синапсы, которые только что выбрали $A$, несут след и ослабевают.

\emph{Неожиданность.} Один провал --- обычная случайность. Но провалы повторяются, и ошибки становятся больше обычного. По гипотезе главы~\ref{sec:nora}, это и есть сигнал \nt{НОРА}: мир изменился. Усиление $g$ падает, выбор становится мягким, и шум чаще приводит к пробе $B$.

\emph{Запись.} По гипотезе главы~\ref{sec:acet}, после перемены \nt{АЦЕТ} поднимается и переводит сеть в режим записи: внутренние связи, хранящие старую привычку, ослаблены, вход от датчиков усилен, веса меняются легко.

\emph{Новая привычка.} Проба $B$ приносит сок, которого никто не ждал: всплеск, $\delta>0$, и синапсы, выбравшие $B$, усиливаются. Несколько таких попыток --- и ожидание $V$ догоняет новый мир, ошибки снова малы.

\emph{Возврат.} Неожиданность прошла, \nt{НОРА} возвращает высокое усиление, выбор снова жёсткий; \nt{АЦЕТ} опускается, сеть в режиме чтения пользуется выученным. Всё это время \nt{СЕРО}, по гипотезе, задаёт горизонт $\tau_\gamma$: насколько далёкий сок ещё стоит ждать.

\begin{figure}[t]
\centering
\begin{tikzpicture}[>=Stealth,x=1.1cm,y=0.95cm]
\foreach \y/\lab in {4/{$\delta$ (\nt{ДОФА})},3/{$g$ (\nt{НОРА})},2/{$a$ (\nt{АЦЕТ})},1/{выбор $B$}}{
  \draw[gray!60!black] (0,\y-0.35) -- (10,\y-0.35);
  \node[left,font=\scriptsize] at (0,\y) {\lab};}
\draw[densely dashed,gray!60!black] (2,0.4) -- (2,4.55) node[above,font=\scriptsize,black] {мир изменился};
\draw[densely dashed,gray!60!black] (5,0.4) -- (5,4.55) node[above,font=\scriptsize,black] {первый сок за $B$};
\pic[scale=0.85] at (2,5.25) {nojuice};
\pic[scale=0.85] at (5,5.25) {juice};
% delta: dips after the change, burst at the first juice for B, then back to zero
\draw[thick,red!70!black] (0,3.65) -- (2.3,3.65) -- (2.3,3.3) -- (2.5,3.3) -- (2.5,3.65) -- (3.2,3.65) -- (3.2,3.3) -- (3.4,3.3) -- (3.4,3.65) -- (4.1,3.65) -- (4.1,3.35) -- (4.3,3.35) -- (4.3,3.65) -- (5.0,3.65) -- (5.0,4.35) -- (5.2,4.35) -- (5.2,3.65) -- (5.8,3.65) -- (5.8,4.1) -- (6.0,4.1) -- (6.0,3.65) -- (6.6,3.65) -- (6.6,3.85) -- (6.8,3.85) -- (6.8,3.65) -- (10,3.65);
% gain: high, drops after surprise, recovers
\draw[thick,blue!60!black] (0,3.2) -- (3.0,3.2) -- (3.4,2.75) -- (5.8,2.75) -- (7.2,3.2) -- (10,3.2);
% ACh: low, rises after the change, falls after learning
\draw[thick,green!45!black] (0,1.75) -- (2.8,1.75) -- (3.3,2.25) -- (6.8,2.25) -- (7.8,1.75) -- (10,1.75);
% choice of B
\draw[thick,black] (0,0.7) -- (3.0,0.7) plot[domain=3:10,samples=40] (\x,{0.7+0.6/(1+exp(-(\x-5.3)*1.8))});
\draw[->,gray!70!black] (0,0.2) -- (10.2,0.2) node[below left,font=\scriptsize,black] {время};
\end{tikzpicture}
\caption{Одно событие, прослеженное через всю машину. Это качественная схема по моделям и гипотезам глав~\ref{sec:dofa}--\ref{sec:acet}, а не результат расчёта. После перемены \nt{ДОФА} отвечает провалами; повторные провалы, по гипотезе, поднимают сигнал неожиданности, \nt{НОРА} снижает усиление, \nt{АЦЕТ} включает запись; первый сок за $B$ даёт всплеск, выбор переходит на $B$, и ручки возвращаются.}
\label{fig:timeline}
\end{figure}

Заметьте, что ни одна ручка не знает, что делают другие. Каждая откликается на свои признаки --- ошибку, её необычную величину, неопределённость, --- и порядок действий складывается сам, как в асинхронной схеме, где каждый блок срабатывает, когда готовы его входы. Такую согласованную работу целиком, насколько нам известно, пока не проверяли: по отдельности ручки измерены, а их совместная работа --- гипотеза.

\section{Чем эта машина не похожа на компьютер}

\begin{table}[t]
\centering
\footnotesize
\setlength{\tabcolsep}{4pt}
\renewcommand{\arraystretch}{1.15}
\begin{tabular}{>{\raggedright}p{2.2cm}>{\raggedright}p{4.2cm}>{\raggedright\arraybackslash}p{6.6cm}}
\toprule
& Обычный компьютер & Мозг \\
\midrule
время & общий такт, около гигагерца & часов нет; каждый нейрон меняется сам, окна задают события и местные ритмы (гл.~\ref{sec:glutcomputer}, \ref{sec:gaba}, \ref{sec:code}) \\
элементы & надёжные двоичные вентили & аналоговые пороговые элементы; синапс теряет большую часть импульсов (гл.~\ref{sec:code}) \\
знак связи & любой & задаётся нейроном-отправителем (гл.~\ref{sec:neurontypes}) \\
память & отдельно от процессора & в тех же связях, что и вычисление, и в самоподдерживающейся активности (гл.~\ref{sec:glutcomputer}, \ref{sec:acet}) \\
обучение & обратное распространение ошибки & местные правила, одно широковещательное число $\delta$ (гл.~\ref{sec:dofa}) и провода ошибки к выходам цепи движений (гл.~\ref{sec:motorlearning}) \\
управление & регистры и шина команд & четыре широковещательных канала, каждый --- жгут из нескольких линий (гл.~\ref{sec:serocomputer}, \ref{sec:dofa}--\ref{sec:acet}) \\
энергия & десятки--сотни ватт на один процессор & около $20$ ватт на весь мозг, в среднем около импульса в секунду на нейрон (гл.~\ref{sec:code}) \\
\bottomrule
\end{tabular}
\caption{Мозг и обычный компьютер.}
\label{tab:computer}
\end{table}

Таблица~\ref{tab:computer} сводит главные отличия. Каждое из них --- не недостаток, который природа не успела исправить, а часть одного решения. Без общего такта нужны датчики «включения---выключения» и детекторы совпадений. Ненадёжным элементам нужна избыточность многих нейронов. Памяти, совмещённой с вычислением, нужен \nt{АЦЕТ}, чтобы запись не портила чтение. Обучению без обратных проводов нужен \nt{ДОФА} и шум. А энергетический предел в один импульс в секунду заставляет весь язык быть скупым и тратить импульсы на новости.

\section{Чего мы не знаем}

Честнее всего закончить эту сводку списком того, чего мы не знаем. Каждый пункт --- задача для инженера.

\emph{Код.} Мы знаем ёмкость импульсного канала и знаем, что получатель выбирает, какую часть сообщения читать (глава~\ref{sec:code}). Но насколько кора пользуется точным временем импульсов и есть ли у нейронов «слова», неизвестно.

\emph{Обучение через много слоёв.} Широковещательный сигнал учит медленно, и медленнее с каждым нейроном, влияющим на результат (утверждение~\ref{prop:broadcastcost}). Как тогда кора настраивает миллиарды весов в многослойных цепях, неизвестно; есть модели, в которых ошибку между слоями разносят пачки импульсов~\cite{Payeur2021}. Возможно, часть ответа --- в вычислениях внутри ветвей самого нейрона, где один нейрон ведёт себя как маленькая многослойная сеть: чтобы повторить ответ модели одного нейрона коры, искусственной сети нужно пять--восемь слоёв~\cite{Beniaguev2021}, а ветви нейронов человека умеют вычислять даже исключающее ИЛИ~\cite{Gidon2020}. Другой кандидат --- самообучение: если кора учится предсказывать собственные входы, ошибка возникает на месте, в каждом слое, и общий сигнал для этого не нужен (глава~\ref{sec:dofa})~\cite{Keller2018}.

\emph{Что на самом деле передают широковещательные каналы.} Для \nt{ДОФА} есть стандартная картина и шесть её расширений (глава~\ref{sec:dofaalt}); для \nt{НОРА} и \nt{АЦЕТ} --- правдоподобные гипотезы (главы~\ref{sec:nora}, \ref{sec:acet}); для \nt{СЕРО} --- в основном догадки.

\emph{Согласование во времени.} Мы видели, как вычислить функцию, отмерить время от события и нарезать поток на окна местным ритмом. Но как огромная сеть без общего такта согласует работу далёких областей, понятно лишь отчасти.

\emph{Энергия.} Двадцать ватт на всё. Мы понимаем, почему код должен быть разреженным (утверждение~\ref{prop:economy}), но построить машину, которая делала бы то же самое при такой же мощности, пока не умеет никто~\cite{MehonicKenyon2022}.

Мозг --- вычислительная машина, собранная не инженером, но по законам, которые любой инженер узнает: RC-цепи, пороги, обратные связи, фильтры, шины и широковещательные регистры. Её схему мы прочитали лишь частично. Дочитывать её будут те, кто умеет читать схемы.

\section{Что дальше в книге}

Эта глава собрала вычислительное ядро машины. Остальные главы выходят за его пределы. Главы~\ref{sec:sensors} и~\ref{sec:recognition} --- о входах: как датчики превращают свет, звук и молекулы в импульсы и как машина узнаёт в них предметы. Главы~\ref{sec:muscles}--\ref{sec:chemoutput} --- о выходах и о том, что их обслуживает: мышцы, боль как сигнал тревоги, обучение движениям по отдельным проводам ошибки и медленный химический выход в тело. Глава~\ref{sec:debugging} --- о том, как машину отлаживают снаружи, в том числе интерфейсами мозг---компьютер. Главы~\ref{sec:eetypes} и~\ref{sec:dendrites} ещё раз сортируют нейроны, теперь по электрической функции и по числу входов. Глава~\ref{sec:sleep} --- о том, что машина делает, когда отключена от мира, а главы~\ref{sec:livingmachine} и~\ref{sec:dishbrain} --- о машинах, собранных из живых нейронов на чипе.

\section{Задачи}

Задачи этой главы соединяют результаты разных глав: каждая опирается по меньшей мере на две из них.

\begin{exercise}[\lvlA]
\label{ex:10:1}
\emph{Шина и регистры.} Сравним потоки данных и управления в машине~\eqref{eq:machine}.
\par\noindent(а)~Оцените сверху поток данных по шине \nt{ГЛУТ}: все $8.6\cdot10^{10}$ нейронов со средней частотой~\eqref{eq:energy}, прочитанные с точностью $1\ms$ по~\eqref{eq:spikeentropy}.
\par\noindent(б)~Оцените поток управления: четыре широковещательных канала, каждый --- жгут примерно из пяти линий (утверждение~\ref{prop:bundle}); значение каждой линии меняется за секунду и различимо на четырёх уровнях. Во сколько раз поток данных больше?
\par\noindent(в)~Сколько времени понадобилось бы каналам управления, чтобы передать столько, сколько шина данных передаёт за одну секунду?
\par\noindent(г)~Пусть широковещательные нейроны \nt{ДОФА}, \nt{СЕРО}, \nt{НОРА} и \nt{АЦЕТ} (таблица~\ref{tab:types}) работают с частотой $5$ импульсов в секунду, а их импульс стоит в $10$--$100$ раз дороже обычного, потому что выходов у них больше. Какую мощность и какую долю энергии сигналов~\eqref{eq:energy} забирает управление? С каким бытовым прибором её можно сравнить?
\end{exercise}

\begin{exercise}[\lvlB]
\label{ex:10:2}
\emph{Две ручки на одной петле.} В~\eqref{eq:machine} усиление $g$ и уровень $a$ умножают одну и ту же скобку. Возьмите петлю \nt{ГЛУТ}--\nt{ГАМК}~\eqref{eq:ei} и поставьте ручки так, как в~\eqref{eq:machine}: $\tau_E\,\dd E/\dd t=-E+g\bigl[(1-a)\,w_{EE}E-w_{EI}I+h\bigr]$, $\tau_I\,\dd I/\dd t=-I+w_{IE}E$; \nt{ГАМК}-нейроны ручками не управляются, передаточная характеристика в рабочей области линейна.
\par\noindent(а)~Выведите, во что превращаются условия утверждения~\ref{prop:ei}.
\par\noindent(б)~Числа рис.~\ref{fig:ei}: $w_{EE}=1.5$, $w_{EI}w_{IE}=2$, $\tau_E=10\ms$, $\tau_I=2\ms$. При $a=0$ найдите усиление $g$, при котором петля начинает колебаться, и частоту колебаний на этой границе.
\par\noindent(в)~Пачка \nt{НОРА} в момент решения поднимает $g$ до $6$. При каком наименьшем $a$ петля остаётся устойчивой? Что будет при $a=0$?
\par\noindent(г)~Что это значит для последовательности на рис.~\ref{fig:timeline}: можно ли крутить ручки \nt{НОРА} и \nt{АЦЕТ} независимо?
\end{exercise}

\begin{exercise}[\lvlB]
\label{ex:10:3}
\emph{Откуда берётся цена широковещания.} Выведите утверждение~\ref{prop:broadcastcost} из~\eqref{eq:perturbation}. Выходы $N$ нейронов $y_k=f(u_k)+\xi_k$ с независимыми гауссовыми $\xi_k$ дисперсии $\sigma^2$ влияют на награду; при малом шуме $R\approx R_0+\sum_kR'_k\xi_k$, и \nt{ДОФА} рассылает $\delta=R-R_0$.
\par\noindent(а)~Покажите, что среднее поправки одного синапса $\delta\,\xi_jx_j$ равно $\sigma^2R'_jx_j$.
\par\noindent(б)~Вычислите её дисперсию.
\par\noindent(в)~Пусть все $R'_k$ одинаковы. Найдите отношение сигнала к шуму одной попытки и число попыток, нужное, чтобы узнать знак поправки с заданной надёжностью. Как оно зависит от $N$?
\par\noindent(г)~Проверьте (в) моделированием на Python при $N=1$, $4$, $16$, $64$ ($5\cdot10^4$ попыток).
\end{exercise}

\begin{exercise}[\lvlB]
\label{ex:10:4}
\emph{Проектирование: связать звук и вспышку.} Событие одновременно звучит и вспыхивает. Звуковой вход приходит на шину через $L_a=5\ms$. Световой --- через $L_v=L_0+t_1$, где $L_0=30\ms$, а $t_1$ --- задержка первого импульса~\eqref{eq:latency} при $\tau=10\ms$ и контрасте от $U=1.2\theta$ до $U=4\theta$. На каждом входе есть и фоновые импульсы, независимые, по $5$ в секунду. Постройте из \nt{ГЛУТ}-нейронов детектор «звук и вспышка от одного события»: линия задержки $d$ на звуковом входе и детектор совпадений~\eqref{eq:window} с порогом $\theta_D$, весом $w$ и постоянной времени $\tau_D$.
\par\noindent(а)~Выберите $d$ и окно так, чтобы связывалось любое событие во всём диапазоне контрастов, а окно было наименьшим.
\par\noindent(б)~Выберите $\tau_D$ при $\theta_D=1.5w$.
\par\noindent(в)~Найдите частоту ложных связываний от фоновых импульсов. Как она зависит от фоновых частот и от разброса задержек? Что это говорит о пользе утверждения~\ref{prop:economy} и о задаче согласования во времени из этой главы?
\end{exercise}

\begin{exercise}[\lvlB]
\label{ex:10:5}
\emph{Моделирование: кто замечает перемену и кто пишет.} Проверим разделение труда из главы~\ref{sec:acet}. Критик следит за ожиданием $s$ по попыткам: в каждой попытке он видит $y=s+\text{шум}$ с дисперсией $R=1$; с вероятностью $h=1/200$ перед попыткой $s$ скачком принимает новое значение из нормального распределения с дисперсией $J^2=4$. Ошибка критика $\delta=y-\hat s$ --- та же ошибка предсказания, что в~\eqref{eq:ctd}. Напишите на Python три критика и для каждого измерьте среднеквадратичную ошибку $\langle(s-\hat s)^2\rangle$ за $10^5$ попыток (десять прогонов).
\par\noindent(а)~Постоянная скорость: $\hat s\leftarrow\hat s+\alpha\delta$. Найдите лучшее $\alpha$.
\par\noindent(б)~Только \nt{АЦЕТ}: фильтр Калмана в попытках, $P\leftarrow P+q$, $K=P/(P+R)$, $\hat s\leftarrow\hat s+K\delta$, $P\leftarrow(1-K)P$, с лучшим постоянным $q$. Объясните результат утверждением~\ref{prop:kalman}.
\par\noindent(в)~\nt{НОРА} и \nt{АЦЕТ}: тот же фильтр при $q=0$, но с детектором неожиданности. Он копит $E\leftarrow\lambda E+\delta^2/(P+R)$ и, когда $E$ превышает порог $c$, поднимает $P$ до $J^2$ (сеть переходит в режим записи) и обнуляет $E$. Подберите $\lambda$ и $c$.
\par\noindent(г)~Найдите нижнюю границу: критик, который точно знает моменты скачков. Что мешает критику (в) до неё дотянуться?
\end{exercise}

\begin{exercise}[\lvlA]
\label{ex:10:6}
\emph{Сколько попыток на смену привычки.} Сеть выучила ожидания $Q_A=0.8$, $Q_B=0.2$ и выбирает по~\eqref{eq:softmax} с $\sigma=1$. Мир меняется: теперь $A$ приносит сок с вероятностью $0.2$, $B$ --- с вероятностью $0.8$. Пока сеть выбирает $A$, ожидание обновляется как $Q_A\leftarrow Q_A+\alpha\,(r-Q_A)$; $B$ почти не пробуется, и $Q_B$ остаётся прежним.
\par\noindent(а)~Найдите среднее $Q_A$ после $n$ выборов $A$.
\par\noindent(б)~При $g=8$ найдите, после скольких выборов $A$ вероятность выбрать $A$ упадёт до $0.6$, при $\alpha=0.1$ и при $\alpha=0.5$ (\nt{АЦЕТ} поднял скорость обучения, глава~\ref{sec:acet}).
\par\noindent(в)~То же, если \nt{НОРА} одновременно опустила $g$ до $2$.
\par\noindent(г)~Почему без проб $B$ вероятность выбрать $A$ никогда не опустится ниже половины, и какая ручка на рис.~\ref{fig:timeline} за это отвечает?
\end{exercise}

\begin{exercise}[\lvlC]
\label{ex:10:7}
\emph{Жгут вместо провода: сколько линий нужно.} В разделе «Чего мы не знаем» остался вопрос, как настраиваются многослойные цепи. Посмотрим, спасает ли дело жгут из $K$ линий (утверждение~\ref{prop:bundle}).
\par\noindent(а)~Модель: $N$ выходов $y_k=w_k+\xi_k$, цель --- $y_k=0$, награда $R=-\sum_ky_k^2$, $w_k(0)=1$. Нейроны разбиты на $K$ групп по $G=N/K$; линия $l$ рассылает своей группе ошибку $\delta_l$ только её выходов, при малом шуме $\delta_l\approx-2\sum_{k\in l}w_k\xi_k$, и $w_k\leftarrow w_k+\eta\,\delta_l\xi_k$. Выведите рекуррентное соотношение для среднего $S=\sum_{k\in l}w_k^2$, найдите лучшее $\eta$ и покажите, что число попыток до $S=\varepsilon S(0)$ равно примерно $(G+2)\ln(1/\varepsilon)$.
\par\noindent(б)~Проверьте моделированием при $N=8$, $32$, $128$ и $K=1$, $4$, $N$, $\sigma=0.1$, $\varepsilon=0.01$.
\par\noindent(в)~Оцените по (а) время обучения цепи из $10^6$ нейронов при жгуте из пяти линий и одной попытке за три секунды.
\par\noindent(г)~Открытый вопрос. Каждой линии жгута нужна своя группа получателей --- это уже адресная, а не широковещательная передача. Сколько линий нужно, чтобы время обучения не зависело от $N$, и откуда могли бы взяться столь многие сигналы ошибки? Сравните с кандидатами из раздела «Чего мы не знаем» и предложите опыт или модель, которые различили бы их.
\end{exercise}
